\id{МРНТИ 61.13.21, 50.43.15}{}

\begin{header}
\swa{К.~Акишев, А.~Зейнуллин, Е.~Айбульдинов, А.~Подвалов, А.~Нурмаганбетов, М.~Чагай, М.~Оспанова, М.~Айтиева, Л.~Мылтыкбаева}{АВТОМАТИЗИРОВАННОЕ ОПРЕДЕЛЕНИЕ СОДЕРЖАНИЯ ХИМИЧЕСКИХ ЭЛЕМЕНТОВ В РУДЕ В УСЛОВИЯХ ФАБРИКИ}

\tsp{1}К.~Акишев\envelope,
\tsp{1}А.~Зейнуллин,
\tsp{1}Е.~Айбульдинов\envelope,
\tsp{2}А.~Подвалов,
\tsp{2}А.~Нурмаганбетов,
\tsp{2}М.~Чагай,
\tsp{1}М.~Оспанова,
\tsp{2}М.~Айтиева,
\tsp{1}Л.~Мылтыкбаева
\end{header}

\begin{affil}
\tsp{1} Казахский университет технологии и бизнеса им. К.Кулажанова, Астана, Казахстан,

\tsp{2}ТОО AG Tech, Астана, Казахстан

\corrauthor{Корреспондент-автор: akmail04cx@mail.ru, elaman\_@mail.ru}
\end{affil}

Актуальность исследования обусловлена необходимостью оперативного и
достоверного определения химического состава руды в условиях
обогатительных фабрик для повышения эффективности рудоподготовки и
снижения энергетических затрат. Целью исследования является разработка и
валидация метода автоматизированного онлайн-анализа содержания меди в
кусковой руде с использованием рентгенорадиометрической технологии. В
качестве основного аналитического признака использован параметр
спектральных отношений H=N\tsb{Cu}/N\tsb{S},
обеспечивающий компенсацию матричных и геометрических эффектов при
поточном измерении.

Методы исследования включали лабораторные и динамические испытания
рентгенорадиометрической сепарации, виртуальную и экспериментальную
сортировку кусковой руды, а также сопоставление аналитических параметров
с результатами химического анализа. Установлена устойчивая линейная
корреляция между параметром H и содержанием меди, что подтверждает
количественную применимость метода. Определён оптимальный диапазон
порога отсечки H ≈ 0,14--0,15, при котором формируется промпродукт с
содержанием Cu 3,0--5,0 \% при выходе 2--12 \% и удалении до 88--98 \%
массы в виде бедных хвостов.

Научная новизна работы заключается в экспериментальном подтверждении
применимости метода спектральных отношений для онлайн-анализа кусковой
руды в фабричном потоке и интеграции аналитического признака в
архитектуру программно-аппаратного комплекса.

Практическая значимость заключается в возможности использования
полученных данных для мониторинга качества рудного потока и оптимизации
режимов предварительного обогащения.

Дальнейшие исследования направлены на расширение спектра определяемых
элементов, опытно-промышленную апробацию и интеграцию аналитических
данных в цифровые модели процессов обогащения.

{\bfseries Ключевые слова:} рентгенорадиометрический анализ, элементный
состав, спектральные отношения, предварительное обогащение, кусковая
руда, программно-аппаратный комплекс.

\begin{header}
БАЙЫТУ ФАБРИКАСЫНДА ФЛОТАЦИЯ ПРОЦЕСІН БАСҚАРУДЫҢ БАҒДАРЛАМАЛЫҚ-АППАРАТТЫҚ КЕШЕНІН ӘЗІРЛЕУ

\tsp{1}К.~Акишев\envelope,
\tsp{1}А.~Зейнуллин,
\tsp{1}Е.~Айбульдинов\envelope,
\tsp{2}А.~Подвалов,
\tsp{2}А.~Нұрмағанбетов,
\tsp{2}М.~Чагай,
\tsp{1}М.~Оспанова,
\tsp{2}М.~Айтиева,
\tsp{1}Л.~Мылтықбаева
\end{header}

\begin{affil}
\tsp{1} Қ. Құлажанов атындағы Қазақ технология және бизнес университеті, Астана, Қазақстан,

\tsp{2}ЖКС AG Tech, Астана, Қазақстан,

e-mail:akmail04cx@mail.ru, elaman\_@mail.ru
\end{affil}

Зерттеудің өзектілігі Кен дайындау тиімділігін арттыру және энергия
шығындарын азайту үшін байыту фабрикалары жағдайында кеннің химиялық
құрамын жедел және сенімді анықтау қажеттілігіне байланысты. Жұмыстың
мақсаты рентген-радиометриялық технологияны қолдана отырып, кесек
Кендегі мыс құрамын автоматтандырылған онлайн-талдау әдісін әзірлеу және
валидациялау болып табылады. Негізгі аналитикалық белгі ретінде h=NCU/NS
спектрлік қатынас параметрі қолданылады, ол ағындық өлшеу кезінде
матрицалық және геометриялық эффектілердің өтемақысын қамтамасыз етеді.

Зерттеу әдістеріне зертханалық және динамикалық рентгендік
радиометриялық сепарация сынақтары, кесек кенді виртуалды және
эксперименттік сұрыптау және аналитикалық параметрлерді химиялық талдау
нәтижелерімен салыстыру кірді. H параметрі мен мыс құрамы арасында
тұрақты сызықтық корреляция орнатылды, бұл әдістің сандық қолданылуын
растайды. H ≈ 0,14-0,15 кесу шегінің оңтайлы диапазоны анықталды, онда
2-12\% шығу және кедей құйрықтар түрінде массаның 88-98\% дейін алып
тастау кезінде cu 3,0--5,0\% өндірістік өнім қалыптасады.

Жұмыстың ғылыми жаңалығы зауыттық ағындағы кесек кенді онлайн талдау
үшін спектрлік қатынас әдісінің қолданылуын эксперименттік растау және
аналитикалық белгіні бағдарламалық-аппараттық кешен архитектурасына
біріктіру болып табылады.

Практикалық маңыздылығы-алынған деректерді кен ағынының сапасын бақылау
және алдын-ала байыту режимдерін оңтайландыру үшін пайдалану мүмкіндігі.

Әрі қарайғы зерттеулер анықталатын элементтердің спектрін кеңейтуге,
тәжірибелік-өнеркәсіптік сынақтан өткізуге және аналитикалық деректерді
байыту процестерінің цифрлық модельдеріне біріктіруге бағытталған.

{\bfseries Түйін сөздер:} рентгендік радиометриялық талдау, элементтік
құрам, спектрлік қатынастар, алдын ала байыту, кесек кен,
бағдарламалық-аппараттық кешен.

\begin{header}
AUTOMATED DETERMINATION OF THE CONTENT OF CHEMICAL ELEMENTS IN ORE IN FACTORY CONDITIONS

\tsp{1}K.~Akishev\envelope,
\tsp{1}А.~Zeynullin,
\tsp{1}Y.~Aibuldinov\envelope,
\tsp{2}A.~Podvalov,
\tsp{2}A.~Nurmaganbetov,
\tsp{2}M.~Chagay,
\tsp{1}M.~Ospanova,
\tsp{2}M.~Aitiyeva,
\tsp{1}L.~Myltykbayeva
\end{header}

\begin{affil}
\tsp{1}Kazakh University of Technology and Business named after K.Kulazhanov, Astana, Kazakhstan,

\tsp{2}AG Tech LLP, Astana, Kazakhstan,

e-mail:akmail04cx@mail.ru, elaman\_@mail.ru
\end{affil}

The relevance of the study is due to the need for prompt and reliable
determination of the chemical composition of ore in processing plants in
order to increase the efficiency of ore preparation and reduce energy
costs. The aim of the study is to develop and validate a method for
automated online analysis of the copper content in lump ore using X-ray
radiometric technology. The spectral ratio parameter H=NCu/NS is used as
the main analytical feature, which provides compensation for matrix and
geometric effects during in-line measurement.

The research methods included laboratory and dynamic tests of X-ray
radiometric separation, virtual and experimental sorting of lump ore, as
well as comparison of analytical parameters with the results of chemical
analysis. A stable linear correlation has been established between the H
parameter and the copper content, which confirms the quantitative
applicability of the method. The optimal range of the cut--off threshold
H ≈ 0.14--0.15 has been determined, at which an industrial product with
a Cu content of 3.0-5.0\% is formed at a yield of 2-12\% and removal of
up to 88-98\% of the mass in the form of poor tailings.

The scientific novelty of the work lies in the experimental confirmation
of the applicability of the spectral relations method for online
analysis of lump ore in a factory stream and the integration of the
analytical feature into the architecture of the software and hardware
complex.

The practical significance lies in the possibility of using the data
obtained to monitor the quality of the ore flow and optimize the
pre-processing modes.

Further research is aimed at expanding the range of identified elements,
pilot testing and integration of analytical data into digital models of
enrichment processes.

{\bfseries Keywords:} X-ray radiometric analysis, elemental composition,
spectral ratios, pre-enrichment, lump ore, hardware and software complex

\begin{multicols}{2}
{\bfseries Введение.} Эффективность процессов дробления, сортировки и
флотационного обогащения на горно-обогатительных предприятиях
определяется точностью и оперативностью информации о химическом, фазовом
составе перерабатываемого сырья.

В условиях изменения минералогического и гранулометрического состава руд
традиционные методы лабораторного химического анализа не всегда,
обеспечивают требуемой оперативности для принятия управленческих решений
в режиме реального времени.

Продолжительность времени между отбором проб и получением результатов
анализа приводят к запаздывающей коррекции технологических режимов,
перерасходу реагентов, снижению извлечения полезных компонентов, росту
энергозатрат.

В настоящее время существующие системы аналитического контроля на
обогатительных фабриках, как правило ориентированы на периодический
лабораторный анализ без интеграции в контуры АСУ-ТП. Современные тренды
отмечают перспективность применения онлайн-аналитических методов, в
частности рентгенофлуоресцентную (XRF), рентгенодифракционную (XRD) и
лазерно-индуцированную спектроскопию (LIBS), используемых в составе
программно-аппаратных комплексов, интегрированных в АСУ ТП предприятия

На практике реализация таких систем сталкивается с рядом технических и
иных ограничений (влияние матричных эффектов, влажности, пылевыделения,
нестабильность гранулометрического состава потока руды и т. д).

Существует проблема пространственно-временной неоднородности состава
руды в потоке-это требует перехода от точечных измерений к непрерывному
мониторингу, интеллектуальной обработке больших массивов спектральных,
технологических и дополнительных данных.

Существующие подходы обработки XRF-спектров, на классических моделях
фундаментальных параметров, снижают точность при изменении матрицы,
толщины слоя материала, ограничивая применение в промышленных потоках.

В этой связи актуальной является разработка программно-аппаратного
комплекса, позволяющего автоматизировать определение содержания
химических элементов и оксидов в руде непосредственно в условиях
фабрики, с возможностью интеграции в АСУ-ТП.

Для решения данной задачи предполагается использование мультисенсорных
аналитических модулей, интеллектуальных методов обработки спектральных
данных, механизмов адаптивной калибровки-это позволит повысить точность
онлайн-анализа, обеспечить оперативную коррекцию технологических
режимов.

Разработка, а также внедрение автоматизированных систем аналитического
контроля химического состава руды, повышает устойчивость технологических
процессов, снижает потери металлов в горно-металлургическом комплексе
Республики Казахстан.

Целью настоящего исследования является разработка и обоснование
программно-аппаратного комплекса для автоматизированного определения
содержания химических элементов и оксидов в рудном сырье в условиях
обогатительной фабрики с возможностью последующей интеграции в систему
АСУ-ТП и использования результатов анализа для оперативной корректировки
режимов переработки.

Экспресс-определение содержания элементов в потоке руды на фабрике
является ключевым условием перехода от «запаздывающего» лабораторного
контроля к оперативному управлению сортировкой, шихтованием и флотацией.

Практика предприятий показывает, что именно online XRF/EDXRF и
родственные радиометрические подходы обеспечивают точность и скорость,
но требуют определенной методики калибровки, защиты оборудования,
компенсации матричных эффектов и интеграции в АСУ- ТП.

В статье {[}1{]} показан промышленный кейс: мониторинг руды на
конвейерах Балхашского и Карагайлинского концентраторов и главном
конвейере Нурказган, измерения спектров выполнялись с секундной
дискретностью, а расчет содержания по меди/цинку/свинцу/железу - по
агрегированию нескольких измерений. Это важно, потому что фиксирует
ключевое требование фабрики: стабилизация оценки за счет временного
усреднения в потоке при высокой неоднородности материала.

В работе {[}2{]} описана методическая сторона внедрения online
XRF-скрининга по Ag и сопутствующим элементам и подчеркиваются
ограничения: низкие концентрации, требования к метрологии и сложности
матрицы. Это напрямую переносимо на вашу задачу определения элементов в
руде «в условиях фабрики», где матричные эффекты и гранулометрия
критичны.

Даже при хорошей сенсорике основная проблема - переносимость калибровки
между участками рудопотока и сменой матриц- исследования в этой области,
а именно подготовке реперных образцов и методам пробоподготовки для
неразрушающего РФА-анализа, показаны в работе {[}3{]} в том числе, выбор
условий гомогенизации, прессования, подбор эталонных образцов,
сопоставление разрушающих/неразрушающих методов).

Работы {[}4,5{]} показывает, что каждое внедрение online XRF требует
предварительных лабораторных исследований, разработки методики и
калибровки по стандартам, а промышленная пригодность определяется
надежностью работы «месяцами без вмешательства». Это фактически
формирует требования к вашему ПАК: пылевлагозащита, стабильность
геометрии, контроль дрейфа.

Исследования посвященные XRF-данным содержащим информацию по множеству
элементов, для ML-подходов представлены в работе {[}6{]}, описана
зависимость- классификация/оценка по XRF-сенсорным данным. В работе
{[}7{]} XRF-сортировки медных руд использует ML-подходы
(оптимизированные SVM/PSO-SVM)-это позволяет преодолевать ограничения
линейных регрессий для сложных руд.

К примеру, в работе {[}8{]} фиксируется проблема не часто показанных
измерений и в тоже время возможность интеллектуального
восстановления/прогноза качества между измерениями.

Работа {[}9{]} показывает использование LIBS в минералогии/переработке,
но это требует сложной калибровки и устойчивости к флуктуациям сигнала.
Гибридная система (XRF + LIBS) улучшает устойчивость и учет всех
параметров, но измерения усложняются.

В работах {[}10,11{]} показано, что состав руды является входным
параметром для шихтования, реагентного режима, стабилизации показателей
при этом к ключевым трудностям относятся - надежные измерения и
переменность входов, в тоже время модели управления должны учитывать
динамику и возмущения.

Анализ зарубежных и казахстанских исследований показывает, что online
определение химического состава руды с использованием
рентгенофлуоресцентных и рентгенорадиометрических методов является
технически реализуемым и уже применяется на отдельных объектах
горно-металлургического комплекса Республики Казахстан, включая
конвейерные линии предприятий корпорации «Казахмыс» и месторождение
Нурказган.

Эти подходы показали эффективность для оперативного мониторинга и
предварительной сортировки, но в большинстве случаев используются как
изолированные измерительные системы, слабо интегрированные в контуры
управления технологическими процессами.

Ограничениями существующих подходов являются:

- высокая чувствительность к матричным эффектам;

- изменения гранулометрического состава;

- влажность;

- необходимость регулярной перекалибровки при изменении геологических и
иных условий.

Это обуславливает потребность в применении интеллектуальных методов
обработки спектральных данных и адаптивных моделей, для учета нелинейных
зависимостей и динамики рудопотока.

В этой связи дальнейшие исследования целесообразно направить на развитие
программно-аппаратных комплексов, объединяющих мультисенсорные
аналитические модули, интеллектуальную обработку данных, адаптивную
калибровку и интеграцию с SCADA и системами оптимального управления.

Для реализации цели в работе были поставлены следующие основные задачи:

- анализ существующих методов онлайн-аналитического контроля химического
состава руды, включая рентгенофлуоресцентную (XRF),
рентгенодифракционную (XRD) и рентгенорадиометрическую сепарацию (РРС),
с точки зрения применимости в потоке и интеграции в контуры АСУ ТП;

- экспериментальное определение химического состава руды
рентгенорадио-метрическим методом;

- формирование аналитического признака, устойчивого к матричным
эффектам;

- виртуальная и динамическая сортировка кусковой руды;

- оценка эффективности предварительного обогащения;

- лабораторная валидация аналитических параметров;

- проектные решения по ПАК (концепция).

Решение указанных задач направлено на создание технологической основы
для внедрения интеллектуальных систем аналитического контроля,
обеспечивающих повышение эффективности обогащения, снижение потерь
металлов и внедрение цифровых технологий в горно-металлургическом
комплексе Республики Казахстан.

{\bfseries Методы и материалы.} \emph{{\bfseries Объект и условия проведения
исследований.}} Объектом

исследования являлись рудные материалы медно-порфирового месторождения
Нурказган.

Экспериментальные работы проводились в рамках опытно-технологических
испытаний с применением рентгенорадиометрической сепарации (РРС) и
лабораторного рентгенофлуоресцентного анализа, направленных на оценку
возможности предварительного обогащения руды и автоматизированного
определения содержания химических элементов.

Исследования выполнялись в лабораторных и пилотных условиях с
использованием дискретных образцов руды и модельных установок
сортировки, без реализации непрерывного промышленного потока на
конвейерных линиях.

В качестве исходного материала использовались технологические пробы руды
месторождения Нурказган. Для экспериментальной сортировки было отобрано
порядка 160 кусков рудного материала машинных классов, включая фракции
порядка −80+50 мм, что соответствует типовым размерам кусковой руды
перед стадией дробления и измельчения.

Для каждой пробы регистрировались:

- интенсивности характеристических рентгеновских линий меди и железа;

- интенсивность рассеянного первичного излучения (параметр Ns);

- интегральные спектры для последующего расчета аналитических признаков.

Лабораторные данные использовались также для сопоставления с
результатами сортировки и оценки достоверности аналитических параметров.

Метод рентгенорадиометрического анализа и формирования разделительного
признака.

В основе измерений использовался метод регистрации характеристического
рентгеновского излучения элементов при облучении руды рентгеновским
источником. Для повышения устойчивости аналитических оценок к изменению
формы, крупности и матричных эффектов применялась нормализация
интенсивности аналитических линий на сигнал рассеянного излучения
образца Ns.

В качестве основного разделительного признака использовался параметр
(ф-ла 1):

\begin{equation}
H=\frac{I_{сu}}{Ns}
\end{equation}

где \(I_{Cu}\) - интенсивность аналитической
линии меди, \(Ns\) - интенсивность рассеянного излучения.

Такой подход позволяет частично компенсировать влияние толщины слоя,
плотности и формы кусков руды, а также вариации матричного состава, что
критично для задач сортировки и предварительного обогащения.

\emph{{\bfseries Процедура экспериментальной сортировки}.} Эксперименты
выполнялись в два этапа:

\emph{{\bfseries Виртуальная сортировка}} - на основе зарегистрированных
спектров проводился расчет аналитического признака \emph{H} и
моделирование разделения материала при различных пороговых значениях
признака. Это позволяло определить предварительный диапазон оптимальных
порогов разделения.

\emph{{\bfseries Динамическая сортировка}} - осуществлялась сортировка руды
в движущемся режиме с регистрацией кумулятивных спектров и параметров
для каждого куска, после чего материал разделялся на «обогащенный» и
«обедненный» продукты.

Для каждого варианта порога разделения рассчитывались показатели:

- выход продуктов;

- содержание меди;

- коэффициент обогащения;

- контрастность разделения;

- эффективность процесса предварительного обогащения.

Статистическая обработка и оценка эффективности. Для анализа
эффективности сортировки использовались зависимости:

- эффективности обогащения от порогового значения признака H;

- контрастности продуктов от среднего содержания меди;

- распределения аналитического признака по классам материала.

Показатели эффективности анализировались в сравнении с исходным
содержанием меди в руде, что позволяло оценить целесообразность
предварительного выделения бедных фракций до стадии флотации.

Оценка выполнялась на усредненных значениях по сериям испытаний, что
снижает влияние случайных флуктуаций сигналов и обеспечивает
статистическую устойчивость выводов.

Использование результатов в концепции программно-аппаратного комплекса.

Полученные экспериментальные данные рассматривались как основа для
формирования требований к программно-аппаратному комплексу
автоматизированного определения химического состава руды. В рамках ПАК
предполагается использование:

- сенсорных модулей рентгенофлуоресцентного и рентгенорадиометрического
анализа;

- подсистемы сбора и обработки спектральных данных;

- модулей вычисления аналитических признаков и формирования управляющих
сигналов для систем сортировки и подготовки сырья.

{\bfseries Результаты и обсуждение.} В ходе лабораторно-пилотных испытаний
руды месторождения Нурказган для каждого куска материала
регистрировались спектры характеристического рентгеновского излучения с
выделением аналитических линий меди и железа, а также интенсивности
рассеянного первичного излучения. Абсолютные значения интенсивностей
аналитических линий демонстрировали значительный разброс, обусловленный
вариациями формы, размеров и плотности кусков руды, а также
неоднородностью минеральной матрицы. Для повышения устойчивости оценки
содержания меди использовалась нормализация интенсивности аналитической
линии меди на интенсивность рассеянного излучения, в результате чего был
сформирован аналитический признак рис.1.
\end{multicols}

\begin{figs}
    \fig[0.45\textwidth]{c/image8}[Рис.1 - Распределение аналитического признака H по серии
образцов руды]
    \fig[0.45\textwidth]{c/image9}[Рис.2 - Зависимость аналитического признака Н от массовой доли Сu]
\end{figs}

\begin{multicols}{2}
Сопоставление аналитического признака H с результатами лабораторного
химического анализа показало устойчивую положительную корреляцию, что
подтверждает возможность использования рентгенорадиометрических
параметров для предварительной количественной оценки содержания меди.
При этом наилучшая воспроизводимость наблюдалась при использовании
усреднённых значений, полученных по кумулятивным спектрам, что снижает
влияние случайных флуктуаций сигнала.

График рис.2 иллюстрирует линейную тенденцию и разброс измерений,
обусловленный неоднородностью руды.

Из рис.2 видно, что наблюдается устойчивая положительная корреляция
между нормализованным спектральным признаком и содержанием меди, что
подтверждает возможность использования рентгенорадиометрических
параметров для предварительной количественной оценки качества руды.
Разброс точек обусловлен неоднородностью минеральной матрицы,
геометрическими факторами характерными для измерений кусковой руды.

Полученные экспериментальные данные показали, что метод обеспечивает
уверенное различение фракций с пониженным и повышенным содержанием меди,
что является достаточным для задач предварительного обогащения и отсечки
бедных кусков до стадии измельчения, и флотации. Абсолютная точность
определения концентрации ниже, чем у лабораторных методов, однако
воспроизводимость, чувствительность признака \emph{H} позволяют
использовать его в качестве технологического индикатора качества руды.

График рис.2 демонстрирует наличие оптимального диапазона порогов
разделения, при котором достигается максимальная эффективность
предварительного обогащения. Смещение порога в сторону меньших значений,
приводит к увеличению выхода пустой породы, а завышение порога
сопровождается потерями полезного компонента.

Полученные экспериментальные данные показали, что метод обеспечивает
различение фракций с пониженным и повышенным содержанием меди - это
является достаточным для задач предварительного обогащения и отсечки
бедных кусков до стадии измельчения и флотации. Абсолютная точность
определения концентрации ниже, чем у лабораторных методов, но
воспроизводимость, чувствительность признака \emph{H} позволяют
использовать его в качестве технологического индикатора качества руды.

На рис.3 представлена зависимость эффективности предварительной
сортировки руды от порогового значения аналитического признака \emph{H},
используемого для разделения материала на обогащённый и обеднённый
продукты. Кривая носит выраженно нелинейный характер и имеет максимум в
области средних значений порога, что указывает на существование
оптимального диапазона разделительных параметров.

При низких значениях порога \emph{H} в обогащённый продукт попадает
значительная доля бедных по меди кусков, что приводит к снижению
контрастности разделения и уменьшению технологической эффективности. В
этом режиме достигается высокий выход концентратного продукта, однако
его среднее содержание меди остается близким к исходному, что снижает
целесообразность предварительного обогащения. При увеличении порога
\emph{H} наблюдается рост содержания меди в обогащённом продукте,
одновременно уменьшается его выход. При чрезмерно высоких порогах,
возрастает доля потерь полезного компонента с отходами, что приводит к
снижению общей эффективности процесса.
\end{multicols}

\begin{figs}
    \fig[0.4\textwidth]{c/image10}[Рис.3 - Зависимость эффективности сортировки от порога
аналитического признака H]
    \fig[0.55\textwidth]{c/image11}[Рис.4 - Аппаратурный спектр руды с выделением аналитической
линии Cu и зоны рассеянного излучения \emph{Ns}]
\end{figs}

\begin{multicols}{2}
Высокая эффективность соответствует режиму, при котором достигается
оптимальное соотношение между извлечением меди и отсечкой пустой породы.
Полученная зависимость подтверждает необходимость адаптивного выбора
порогов сортировки в зависимости от текущих свойств руды, а также
производственных требований. Результаты демонстрируют, что управление
порогом аналитического признака является эффективным инструментом
оптимизации предварительного обогащения и может использоваться в составе
программно-аппаратного комплекса, как регулируемый параметр
технологического режима.

На рис.4 представлен аппаратурный спектр руды месторождения Нурказган,
зарегистрированный в процессе рентгенорадиометрического анализа
кускового материала. В спектре отчетливо выделяются характеристические
пики меди, а также широкая область рассеянного излучения, формируемая в
результате взаимодействия первичного рентгеновского пучка с минеральной
матрицей образца.

Анализ спектров показал, что абсолютная интенсивность аналитической
линии меди существенно зависит не только от концентрации элемента, но и
от геометрии куска, толщины слоя, плотности материала и минерального
состава. Эти факторы вызывают матричные эффекты и приводят к
значительному разбросу аналитических параметров, что делает
использование ненормализованных интенсивностей недостаточно устойчивым
для задач сортировки и оперативного контроля качества. Для компенсации
указанных эффектов была использована нормализация интенсивности
аналитической линии меди на интенсивность рассеянного излучения,
определяемого в спектре в отдельном энергетическом диапазоне (зона
\emph{Ns} на рис.4).

Применение данного признака позволило существенно снизить влияние
вариаций геометрии кусков и неоднородности минеральной матрицы, параметр
Ns отражает суммарную интенсивность рассеяния, коррелирует с массой и
плотностью материала в зоне облучения. Отношение
I\tsb{Cu}/N\tsb{s} является более
устойчивым к изменению условий измерений по сравнению с использованием
абсолютных интенсивностей.

Экспериментальные распределения значений признака \emph{H} показали
формирование устойчивых групп, по различным уровням содержания меди, что
обеспечивает его информативность и пригодность для порогового разделения
материала при предварительном обогащении. Использование данного признака
обеспечило повышение контрастности разделения, позволило сформировать
стабильные пороговые значения для сортировки в различных сериях
испытаний.

Устойчивость аналитического признака \emph{H} подтверждена для условий
кусковой руды в лабораторно-пилотных экспериментах. Влияние факторов
промышленной среды, таких как высокая влажность, запылённость, работа в
потоках пульпы, определяет необходимость дальнейшей опытно-промышленной
валидации метода.

Таким образом, формирование нормализованного аналитического признака на
основе отношения интенсивности аналитической линии меди к интенсивности
рассеянного излучения является эффективным способом компенсации
матричных эффектов и обеспечивает технологически устойчивый параметр для
использования в составе программно-аппаратного комплекса
автоматизированного контроля качества руды.

В рамках выполненных исследований реализован подход виртуальной
сортировки кусковой руды, при котором решение о принадлежности каждого
куска к обогащённому или отвальному продукту принимается на основе
аналитического признака, сформированного по результатам
рентгенорадиометрического измерения, без фактического механического
разделения материала. Он позволяет оценить потенциальную эффективность
сортировки, оптимизировать пороговые значения признака и спрогнозировать
технологические показатели процесса без проведения дорогостоящих
промышленных испытаний.

На рис.5 представлена схема экспериментальной установки и геометрия
измерений, в которой каждый кусок руды последовательно проходит зону
облучения и регистрации спектра. После вычисления аналитического
признака каждому образцу присваивается виртуальный класс ---
«концентрат» или «хвосты» --- в зависимости от заданного порога H. Далее
на основе лабораторных данных химического анализа рассчитываются
показатели выхода, извлечения и контрастности разделения для различных
пороговых значений.
\end{multicols}

\fig[0.7\textwidth]{c/image12}[Рис.5 - Схема виртуальной сортировки кусковой руды]

\begin{multicols}{2}
Динамический характер сортировки проявляется в том, что параметры
разделения могут изменяться в зависимости от текущих характеристик
поступающей руды, а также при изменении среднего содержания меди и
гранулометрического состава целесообразно адаптировать порог
аналитического признака для сохранения оптимального соотношения между
извлечением полезного компонента и отсечкой пустой породы. Что
соответствует концепции динамической сортировки, при которой управляющий
параметр изменяется во времени в зависимости от входных условий.

График на рис.6 демонстрирует наличие оптимальной зоны порогов, при
которых достигается максимальная эффективность предварительного
обогащения. При снижении порога увеличивается выход, но падает качество
продукта, тогда как при завышении порога возрастает доля потерь
полезного компонента.

Зависимость контрастности разделения и коэффициента обогащения от
среднего содержания меди в исходной руде, подтверждает необходимость
адаптивной настройки параметров сортировки при изменении качества
питания. Эти результаты показывают, что фиксированное пороговое значение
не обеспечивает устойчивых технологических показателей при вариациях
свойств руды.

Виртуальная сортировка позволяет определить потенциальный эффект
предварительного обогащения, сформировать правила управления процессом,
а динамическая сортировка рассматривается как развитие данного подхода,
ориентированное на адаптацию параметров разделения в режиме реального
времени. Полученные результаты подтверждают целесообразность интеграции
модуля сортировки в состав программно-аппаратного комплекса и
использования аналитических признаков в контуре автоматизированного
управления подготовкой сырья перед стадией измельчения и флотации.
\end{multicols}

\begin{figs}
    \fig[0.45\textwidth]{c/image13}[Рис.6 - Зависимость эффективности сортировки от порогового
значения аналитического признака H]
    \fig[0.45\textwidth]{c/image14}[Рис.7 - Принцип рентгенорадиометрической сепарации]
\end{figs}

\begin{multicols}{2}
Продолжение развития метода является переход от виртуальной сортировки к
опытно-промышленным испытаниям с использованием исполнительных
механизмов (воздушные или механические отсекатели) и интеграцией
алгоритмов принятия решений в систему управления технологическим
процессом.

Принцип рентгенорадиометрической сепарации (РРС) заключается в том, что
рентгеновская трубка возбуждает характеристическое рентгеновское
излучение в куске руды. Детектор регистрирует интенсивность линий меди и
рассеянного излучения, после чего вычисляется аналитический параметр H =
N\tsb{C}u/N\tsb{S}. На основе сравнения параметра с
заданным порогом выполняется разделение потока на «проход» (в дальнейшую
переработку) и «отсев» (в отвал или буфер). Использование спектрального
отношения позволяет компенсировать влияние размеров и формы кусков, а
также матричного состава породы.

График на рис.8 демонстрирует устойчивую положительную корреляцию между
аналитическим параметром \emph{H} и фактическим содержанием меди в руде.
Это подтверждает корректность выбора признака сортировки и возможность
использования параметра не только для бинарной отсечки пустой породы, но
и для управления качеством питания обогатительной фабрики. При
увеличении значения \emph{H} наблюдается рост содержания \emph{Cu}, что
позволяет реализовать селективное управление рудопотоком.
\end{multicols}

\begin{figs}
    \fig[0.52\textwidth]{c/image15}[Рис.8 - Зависимость содержания Cu от аналитического признака H]
    \fig[0.4\textwidth]{c/image16}[Рис.9 - Эффективность сортировки от порога отсечки H]
\end{figs}

\begin{multicols}{2}
На графике рис.9 показано, что эффективность предварительного обогащения
имеет выраженный максимум в области \emph{H} ≈ 0,14--0,15. При более
низких значениях порога происходит недостаточное отделение пустой
породы, при более высоких - теряется часть полезного компонента с
отсевом. Выбор оптимального порога является ключевым фактором
технологической эффективности и должен уточняться при
опытно-промышленных испытаниях.

Диаграмма рис.10 иллюстрирует характерную для предварительного
обогащения ситуацию, при которой основная часть массы руды (около
70--90\%) отсекается как бедная и не направляется на измельчение, тогда
как меньшая часть потока, обладающая повышенным содержанием полезного
компонента, поступает в дальнейшую переработку. Это обеспечивает
значительное снижение энергозатрат на измельчение и повышение
эффективности последующих стадий обогащения.
\end{multicols}

\begin{figs}
    \fig[0.40\textwidth]{c/image17}[Рис.10 - Баланс рудопотока после предварительного обогащения]
    \fig[0.55\textwidth]{c/image18}[Рис.11 - Лабораторная зависимость содержания Cu от
аналитического параметра H]
\end{figs}

\begin{multicols}{2}
Целью лабораторных испытаний являлась проверка пригодности
аналитического параметра \emph{H=N\tsb{Cu}/N\tsb{S}}
в качестве устойчивого разделительного признака для сортировки руды
месторождения Нурказган методом рентгенорадиометрической сепарации.

Валидация включала:

- измерение параметра \emph{H} для отдельных кусков руды;

- определение фактического содержания меди лабораторным анализом;

- анализ зависимости между \emph{H} и \emph{Cu;}

- оценку эффективности сортировки при различных порогах отсечки.

На рис.11 представлена зависимость содержания меди от аналитического
параметра \emph{H}, полученная по лабораторным данным.

Из рис.11 видно, что наблюдается устойчивая положительная корреляция
между H и содержанием \emph{Cu.}

При H <{} 0,08 содержание меди, как правило, не превышает
0,5--0,7 \%.

При H ≈ 0,14--0,16 содержание меди возрастает до 1,2--1,5 \%.

При H \textgreater{} 0,2 содержание меди превышает 1,8--2,5 \%.

Это подтверждает, что параметр \emph{H} является информативным и может
использоваться для отсечки пустой породы и селективного управления
качеством рудопотока.

На рис.12 приведена зависимость эффективности предварительного
обогащения от значения порога отсечки аналитического параметра \emph{H}.

Анализ графика показывает:

- при низких порогах (\emph{H} <{} 0,1) эффективность сортировки
низкая из-за слабого отделения пустой породы;

- максимальная эффективность достигается при \emph{H} ≈ 0,14--0,15;

- при дальнейшем увеличении порога часть обогащенной руды начинает
уходить в отсев, что снижает общий эффект.

Таким образом, область \emph{H} ≈ 0,14--0,15 может рассматриваться как
оптимальный стартовый диапазон порога отсечки, подлежащий дальнейшей
оптимизации в опытно-промышленных условиях.
\end{multicols}

\begin{figs}
    \fig[0.45\textwidth]{c/image19}[Рис.12 - Эффективность предварительного обогащения в функции
порога H]
    \fig[0.43\textwidth]{c/image20}[Рис.13 - Классификация зон эффективности сортировки]
\end{figs}

\begin{multicols}{2}
На основании лабораторных данных режимы сортировки могут быть условно
разделены на три зоны рис.13:

- не эффективная - недостаточное отделение пустой породы;

- эффективное-оптимальный баланс между извлечением и отсевом;

- весьма эффективное-максимальное обогащение при ограниченном выходе.

Наибольшая доля наблюдений относится к эффективной зоне, что
подтверждает корректность выбранного аналитического параметра и
принципиальную реализуемость технологии предварительного обогащения для
данной руды.

В таблице 1 приведены результаты лабораторной проверки зависимости
аналитического параметра H=N\tsb{Cu}/N\tsb{S} от
фактического содержания меди и эффективности предварительного
обогащения.
\end{multicols}

\tcap{Таблица 1 - Лабораторная валидация аналитического параметра
\emph{H}}
\begin{longtblr}[
  label = none,
  entry = none,
]{
  width = 0.7\linewidth,
  colspec = {Q[40]Q[317]Q[219]Q[362]},
  cells = {c},
  cell{1}{2} = {font=\bfseries},
  cell{1}{3} = {font=\bfseries},
  cell{1}{4} = {font=\bfseries},
  hlines,
  vlines,
}
№  & Аналитический параметр H & Содержание Cu, \% & Эффективность сортировки, \% \\
1  & 0.03                     & 0.3               & 15                           \\
2  & 0.05                     & 0.45              & 25                           \\
3  & 0.07                     & 0.6               & 40                           \\
4  & 0.09                     & 0.8               & 55                           \\
5  & 0.11                     & 0.95              & 70                           \\
6  & 0.13                     & 1.1               & 82                           \\
7  & 0.145                    & 1.25              & 90                           \\
8  & 0.16                     & 1.45              & 88                           \\
9  & 0.18                     & 1.65              & 80                           \\
10 & 0.22                     & 1.9               & 65                           \\
11 & 0.26                     & 2.2               & 50                           \\
12 & 0.3                      & 2.5               & 40                           
\end{longtblr}

\begin{multicols}{2}
По результатам лабораторных испытаний установлено, что:

- аналитический параметр \emph{H} обладает высокой чувствительностью к
содержанию меди;

- обеспечивается контрастное разделение продуктов;

- существует выраженный оптимум порога сортировки;

- параметр пригоден для внедрения в промышленные системы сортировки и
цифровые контуры управления.

На рис.14 представлена архитектура программно-аппаратного комплекса
предварительного обогащения построена по иерархическому принципу и
включает четыре уровня. На измерительном уровне выполняется возбуждение
характеристического рентгеновского излучения и регистрация спектров от
кусков руды, формируются первичные сигналы по меди и рассеянному
излучению. На уровне контроллера сепаратора осуществляется расчет
аналитического параметра \emph{H}, временная привязка сигнала к
конкретному куску и формирование команды на исполнительные механизмы для
разделения потока.

Уровень PLC/SCADA обеспечивает управление технологическим оборудованием,
синхронизацию работы узла сортировки с конвейерным транспортом, сбор
статистики и передачу данных на верхний уровень.

На уровне MES/AI-модуля выполняется анализ накопленных данных,
оптимизация порога отсечки (cut-off) и интеграция ПАК в общий контур
управления фабрикой и цифровой двойник процесса. Такая структура
обеспечивает как автономную работу узла сортировки, так и его адаптивную
настройку в зависимости от качества руды и технологических показателей
обогащения.
\end{multicols}

\fig[0.7\textwidth]{c/image21}[Рис.14 - Архитектура программно-аппаратного комплекса
предварительного обогащения]

\begin{multicols}{2}
Полученные результаты, показывает, что научная новизна исследования
заключается:

- в экспериментальном подтверждении применимости метода спектральных
отношений для промышленного онлайн-анализа именно кусковой руды, тогда
как большинство известных решений ориентированы на пульпу или
порошкообразные пробы. Впервые показано, что параметр \emph{H} сохраняет
линейную корреляцию с фактическим содержанием меди в широком диапазоне
технологических условий, что позволяет использовать его не только для
мониторинга состава, но и как количественный признак для классификации
рудного потока;

- интеграции аналитического параметра в архитектуру
программно-аппаратного комплекса с возможностью адаптивной корректировки
порога отсечки, что формирует основу для перехода от пассивного
химического контроля к элементам интеллектуального управления качеством
рудопотока.

Выполненное исследование расширяют область применения
рентгенорадиометрических методов от задач контроля к задачам оптимизации
технологических процессов на основе данных химического анализа.

{\bfseries Выводы.} В работе решён комплекс задач, направленных на
разработку метода автоматизированного определения химического состава
руды в условиях обогатительной фабрики. Сравнительный анализ методов
онлайн-контроля (XRF, XRD и РРС) показал, что рентгенорадиометрический
метод является наиболее пригодным для поточного анализа кусковой руды
без пробоподготовки и легко интегрируется в контуры АСУ ТП.

Экспериментально подтверждена возможность определения содержания меди на
основе характеристического рентгеновского излучения с использованием
аналитического параметра H=N\tsb{Cu}/N\tsb{S}
обеспечивающего компенсацию матричных и геометрических эффектов.

Установлена устойчивая линейная корреляция между параметром \emph{H} и
фактическим содержанием \emph{Cu}, что позволяет использовать данный
признак для классификации руды по качеству.

По результатам виртуальной и динамической сортировки определён
оптимальный диапазон порога отсечки \emph{H} ≈ 0,14--0,15. В
лабораторных условиях получено формирование промпродукта с содержанием
\emph{Cu} 3,0--5,0 \% при выходе 2--12 \% и удаление до 88--98 \% массы
в виде бедных хвостов, подтверждающее эффективность предварительного
обогащения.

Разработана концепция программно-аппаратного комплекса, обеспечивающего
автоматизированный сбор, обработку аналитических данных, а также
возможность интеграции в АСУ-ТП фабрики.

Дальнейшие исследования целесообразно направить на опытно-промышленную
валидацию метода, расширение спектра определяемых элементов и интеграцию
аналитических данных в цифровые модели процессов для оптимизации режимов
обогащения.

\emph{{\bfseries Финансирование.}} Статья подготовлена по результатам
исследований в рамках ПЦФ BR28713037 «Разработка программно-аппаратного
комплекса на базе технологий искусственного интеллекта для управления
процессом флотации на горно-обогатительных предприятиях».
\end{multicols}

\begin{center}
{\bfseries Литература}
\end{center}

\begin{refs}
1. Nigmatulin A., Abdrakhmanova Z., Kan A., Efimenko S., Makarov D.
Online X-Ray Fluorescence Monitoring of Coarse Ore for Silver at the
Process Conveyors at Kazakhmys Corporation LLC // Journal of the Polish
Mineral Engineering Society-2020.-Vol.2(1).-P.145-148. DOI
\href{http://doi.org/10.29227/IM-2020-01-56}{10.29227/IM-2020-01-56}.

2. Abdrakhmanova Z., Kan A., Yun R., Yefimenko S., Makarov D. Online
X-ray fluorescence screening of ore for silver at the mines operated by
Kazakhmys corporation LLC// Sustainable Extraction and Processing of Raw
Materials Journal\emph{.-}2021.-Vol.2(2).-P.6-11.
\href{https://doi.org/10.5281/zenodo.5594992}{DOI
10.5281/zenodo.5594992}.

3. Колбин В.В. Методы подготовки реперных образцов почв для
неразрушающего РФА-анализа//Вестник НЯЦ РК\emph{.-} 2022.-Т.3(91).- С.
11-20. DOI 10.52676/1729-7885-2022-3-11-20.

4. Bardzinski P. et al. (2020). Copper Ore Quality Tracking in a Belt
Conveyor System Using simulation tools // Natural Resources
Research.-2019.-Vol.29(11)- P.1031-1040. DOI
\href{https://doi.org/10.1007/s11053-019-09493-6?urlappend=\%3Futm_source\%3Dresearchgate.net\%26utm_medium\%3Darticle}{10.1007/s11053-019-09493-6}.

5. Sokolov A., Kuzmovs V., Ordóñez U.M., Gostilo V. Online XRF Analysis
of Elements in Minerals on a Conveyor Belt // Mining. -2025.-Vol.
5(4):77. \href{https://doi.org/10.3390/mining5040077}{DOI
10.3390/mining5040077}.

6. Xu Y. Development of Methodological Framework for XRF Sensor Based
Ore Sorting Utilizing Machine Learning Approaches// PhD thesis, Mining
Engineering (Vancouver)-2024.- 241p. URL:
\url{https://share.google/rEWpSE5KHuutY6DiO}. Date of access: 26.12.2025

7. Liu Z. et al. Enhancing XRF sensor-based sorting of porphyritic
copper or using particle swarm optimization-support vector machine
(PSO-SVM) algorithm //International Journal of Mining Science and
Technology\emph{. -}2024.-Vol.34(4)\emph{. -}P.545-556. DOI
\href{https://doi.org/10.1016/j.ijmst.2024.04.002}{10.1016/j.ijmst.2024.04.002}.

8. Peng C. et al. Study on Prediction of Zinc Grade by Transformer Model
with De-Stationary Mechanism// Minerals.-2024.-Vol.14(3):230.
\href{https://doi.org/10.3390/min14030230}{DOI 10.3390/min14030230}.

9. Harmon R.S. et al. Laser-Induced Breakdown Spectroscopy in Mineral
Exploration and Ore Processing// Minerals.-2024.-Vol.~14(7): 731.
\href{https://doi.org/10.3390/min14070731}{DOI 10.3390/min14070731}.

10. Oosthuizen D.J., Craig~I.K., Jämsä-Jounela S-L.,
\href{https://www.sciencedirect.com/author/55460745900/bei-sun}{Sun} B.
On the current state of flotation modelling for process
control//\href{https://www.sciencedirect.com/journal/ifac-papersonline}{IFAC-PapersOnLine}.-2017.-Vol.50(2).-P.19-24.
DOI 10.1016/j.ifacol.2017.12.004.

11. Salo M. et al.Active sensing in froth flotation//Int. J. Mining and
Mineral Engineering.-2024.- Vol.15(2). -P.111-130. DOI
\href{https://doi.org/10.1504/IJMME.2024.10063891}{10.1504/IJMME.2024.10063891}.
\end{refs}

\begin{center}
{\bfseries References}
\end{center}

\begin{refs}
1. Nigmatulin A., Abdrakhmanova Z., Kan A., Efimenko S., Makarov D.
Online X-Ray Fluorescence Monitoring of Coarse Ore for Silver at the
Process Conveyors at Kazakhmys Corporation LLC // Journal of the Polish
Mineral Engineering Society-2020.-Vol.2(1).-P.145-148. DOI
\href{http://doi.org/10.29227/IM-2020-01-56}{10.29227/IM-2020-01-56}.

2. Abdrakhmanova Z., Kan A., Yun R., Yefimenko S., Makarov D. Online
X-ray fluorescence screening of ore for silver at the mines operated by
kazakhmys corporation LLC// Sustainable Extraction and Processing of Raw
Materials Journal\emph{.-}2021.-Vol.2(2).-P.6-11. DOI
10.5281/zenodo.5594992.

3. Kolbin V.V. Metody podgotovki repernyh obrazcov pochv dlja
nerazrushajushhego RFA-analiza//Vestnik NJaC RK.- 2022.-T.3(91).- S.
11-20. DOI 10.52676/1729-7885-2022-3-11-20.{[}in Russian{]}

4. Bardzinski P. et al. (2020). Copper Ore Quality Tracking in a Belt
Conveyor System Using simulation tools // Natural Resources
Research.-2019.-Vol.29(11)- P.1031-1040. DOI 10.1007/s11053-019-09493-6.

5. Sokolov A., Kuzmovs V., Ordóñez U.M., Gostilo V. Online XRF Analysis
of Elements in Minerals on a Conveyor Belt // Mining. -2025.-Vol.
5(4):77. \href{https://doi.org/10.3390/mining5040077}{DOI
10.3390/mining5040077}.

6. Xu Y. Development of Methodological Framework for XRF Sensor Based
Ore Sorting Utilizing Machine Learning Approaches// PhD thesis, Mining
Engineering (Vancouver)-2024.- 241p. URL:
\url{https://share.google/rEWpSE5KHuutY6DiO}. Date of access: 26.12.2025

7. Liu Z. et al. Enhancing XRF sensor-based sorting of porphyritic
copper or using particle swarm optimization-support vector machine
(PSO-SVM) algorithm //International Journal of Mining Science and
Technology\emph{.-}2024.-Vol.34(4)\emph{.-}P.545-556. DOI
\href{https://doi.org/10.1016/j.ijmst.2024.04.002}{10.1016/j.ijmst.2024.04.002}.

8. Peng C. et al. Study on Prediction of Zinc Grade by Transformer Model
with De-Stationary Mechanism// Minerals.-2024.-Vol.14(3):230.
\href{https://doi.org/10.3390/min14030230}{DOI 10.3390/min14030230}.

9. Harmon R.S. et al. Laser-Induced Breakdown Spectroscopy in Mineral
Exploration and Ore Processing// Minerals.-2024.-Vol.~14(7): 731.
\href{https://doi.org/10.3390/min14070731}{DOI 10.3390/min14070731}.

10. Oosthuizen D.J., Craig~I.K., Jämsä-Jounela S-L.,
\href{https://www.sciencedirect.com/author/55460745900/bei-sun}{Sun} B.
On the current state of flotation modelling for process
control//\href{https://www.sciencedirect.com/journal/ifac-papersonline}{IFAC-PapersOnLine}.-2017.-Vol.50(2).-P.19-24.DOI
10.1016/j.ifacol.2017.12.004.

11. Salo M. et al. Active sensing in froth flotation//Int. J. Mining and
Mineral Engineering. -2024.- Vol.15(2). -P.111-130. DOI
\href{https://doi.org/10.1504/IJMME.2024.10063891}{10.1504/IJMME.2024.10063891}.
\end{refs}

\begin{info}
\hspace{1em}{\bfseries \emph{Сведения об авторах}}

Акишев К.М. -к.т.н., ассоциированный профессор, Казахский университет
технологии и бизнеса им. К. Кулажанова, Астана, Казахстан, e-mail:
akmail04cx@mail.ru;

Зейнуллин А.А.- д.т.н., профессор, Казахский университет технологии и
бизнеса им. К. Кулажанова, Астана, Казахстан, e-mail: karim\_57@mail.ru;

Айбульдинов Е.К. -- доктор PhD, Казахский университет технологии и
бизнеса им. К. Кулажанова, Астана, Казахстан, e-mail:
elaman\_@mail.ru;

Подвалов А.Ю.- к.г.н, ТОО «AG TECH», генеральный директор, Астана,
Казахстан, e-mail: a.podvalov@ag-tech.kz;

Нурмаганбетов А.С.- к.т.н, доктор философии (PhD), ассоциированный
профессор, ТОО «AG TECH», Астана, Казахстан, e-mail:
a.nurmaganbetov@ag-tech.kz;

Чагай М. А.- ТОО «AG TECH», Астана, Казахстан, e-mail:
m.chagay@ag-tech.kz;

Оспанова М.К. - к.ф.-м.н., Казахский университет технологии и бизнеса
им. К. Кулажанова, Астана, Казахстан, e-mail:
ospanova\_madi\_59@mail.ru;

Айтиева М. И.-ТОО «AG TECH», Астана, Казахстан, e-mail:
m.aitiyeva@ag-tech.kz;

Мылтыкбаева Л.А.- к.т.н., Казахский университет технологии и бизнеса им.
К. Кулажанова, Астана, Казахстан, e-mail: L.multykbayeva@gmail.com.

\hspace{1em}\emph{{\bfseries Information about the authors}}

Akishev K. - Candidate of Technical Sciences, Associate Professor, K.
Kulazhanov Kazakh University of Technology and Business, Astana,
Kazakhstan, e-mail: akmail04cx@mail.ru;

Zeynullin A.- Doctor of Technical Sciences, Professor, K. Kulazhanov
Kazakh University of Technology and Business, Astana, Kazakhstan,
e-mail: karim\_57@mail.ru;

Aibuldinov Y.- PhD, K. Kulazhanov Kazakh University of Technology and
Business, Astana, Kazakhstan, e-mail:
elaman\_@mail.ru;

Podvalov A.- PhD in Geographic sciences, AG TECH LLP, Kazakhstan,
Astana, e-mail:
a.podvalov@ag-tech.kz;

Nurmaganbetov A.- Candidate of Technical Sciences, Associate Professor,
AG TECH LLP, Astana, Kazakhstan, e-mail: a.nurmaganbetov@ag-tech.kz;

Chagai M.- AG TECH LLP, Astana, Kazakhstan, e-mail: m.chagay@ag-tech.kz;

Ospanova M. - Candidate of Physical and Mathematical Sciences, K.
Kulazhanov Kazakh University of Technology and Business, Astana,
Kazakhstan, e-mail: ospanova\_madi\_59@mail.ru;

Aitieva M.- AG TECH LLP, Astana, Kazakhstan, e-mail:
m.aitiyeva@ag-tech.kz;

Myltykbayeva L - Candidate of Technical Sciences, K. Kulazhanov Kazakh
University of Technology and Business, Astana, Kazakhstan, e-mail:
L.multykbayeva@gmail.com.
\end{info}
